\documentclass[12pt]{article}

% Packages 
\usepackage[left=1cm, right=1cm, top=1cm, bottom=1.2cm]{geometry}
\usepackage{fontspec}
\usepackage{xeCJK}
\setCJKmainfont{Noto Serif CJK SC} 

\usepackage{packages}
\usepackage{scribe}
\usepackage{listings}
\usepackage{macros}
\usepackage[
backend=biber,
style=numeric,
citestyle=numeric
]{biblatex}
\addbibresource{references.bib}

% Document Settings
\setlength{\parindent}{2em}
\Scribe{Luke Zhang}
\LectureNumber{N}
\LectureDate{\today}
\LectureTitle{Practical Manifold Muon}

\begin{document}
	\MakeScribeTop

\section{Introduction}

In this article, we will be trying to collect fragmented knowledge about the newly proposed optimizer "Manifold Muon" by Jeremy Bernstein \cite{bernstein2025manifolds} to construct something that doesn't only exist in theory, but practically usable. 

The most intuitive way to understand manifold muon is perhaps through this table: 

\medskip

\begin{center}
\begin{tabular}{|c|c|c|}
\hline
Manifold & Norm & Optimizer \\
\hline \hline
Euclidean $\mathbb{R}^n$ & Euclidean Norm & Vanilla Gradient Descent \\ \hline
Matrix Space $\mathbb{R}^{m \times n}$ & Spectral Norm & Muon \\ \hline
Stiefel Manifold $\subset \mathbb{R}^{m \times n}$ & Spectral Norm & Manifold Muon \\ \hline
\end{tabular}
\end{center}

Seeing the similarity between Muon and Manifold Muon, let's start looking by looking at the derivation of Muon and how it's being practically used right now. As we'll show later, much of the techniques and strategies transfer well to Manifold Muon.

\subsection{Notation}

We use the following notation throughout the article.
Let $\Weights$ denote the weight matrix,
$\Gr$ the gradient matrix,
$\Momentum$ the momentum-augmented gradient matrix,
and $\Orthogonal$ an orthogonalized matrix.


\section{Understanding Muon}
Muon is a recently proposed optimizer that updates the weight space instead of individual vectors. The algorithm written down by Keller Jordan is as follows \cite{jordan2024muon}:

\begin{algorithm}
\caption{Muon}
\begin{algorithmic}[1]
    \Require Learning rate $\eta$, momentum $\mu$
    \State Initialize: $\Momentum[0] \gets 0$
    \For{$t = 1, \dots$}
        \State Compute gradient $\Gr[t] \gets \grad[\Weights]\lossfunc(\Weights[t-1])$
        \State $\Momentum[t] \gets \mu \Momentum[t-1] + \Gr[t]$
        \State $\Orthogonal[t] \gets$ NewtonSchultz5$(\Momentum[t])$
        \State Update parameters $\Weights[t] \gets \Weights[t-1] - \eta\Orthogonal[t]$
    \EndFor

    \State \Return $\Weights[t]$
\end{algorithmic}
\end{algorithm}

Observe that the Muon algorithm is essentially SGD Momentum, but with one extra caveat: the NewtonSchultz5 function on step 5. To understand this function, let's quickly gloss over the derivation of Muon, you can read more about the specifics in Jeremy Bernstein's blog post \cite{bernstein2025deriving}. 

\subsection{Deriving Muon (quickly)}

Ignoring SGD Momentum, the most barebones update rule for Muon is as follows:
\begin{align}
    \Weights \gets \Weights - \eta \times \text{NewtownSchulz}(\grad[\Weights]\lossfunc) \label{base-muon-update}
\end{align}

To get there, let's first define the norm we'll be using: the spectral norm over the weight space, defined as:
\begin{align}
    \snorm{\Weights} := \max_{x \ne 0} \frac{\snorm{\Weights x}}{\snorm{x}}
\end{align}

We now want to find a weight update that maximizes improvement in loss $\lossfunc$ while having some constraint on the change in the output ($\change y$) we want to make in one step. In other words, we are solving the following problem:
\begin{align}
    \min_{\change \Weights} \langle \grad[\Weights]\lossfunc, \change \Weights \rangle \text{ subject to } \snorm{\change y} \le \eta \label{muon}
\end{align}

Starting from the definition of $\change y$: 
\begin{align*}
    \change y &= (\Weights + \change \Weights)x - \Weights x \\
    \change y &= \change \Weights x \\
    \snorm{\change y} &= \snorm{\change \Weights x} \le \snorm{\change \Weights} \times \snorm{x}
\end{align*}

Substitute into \probref{muon}:
\begin{align}
    \min_{\change \Weights} \langle \grad[\Weights]\lossfunc, \change \Weights \rangle \text{ subject to } \snorm{\change \Weights} \le \eta \label{muon-final}
\end{align}

Notice that $\eta$ in \probref{muon-final} should be scaled by a factor of $\frac{1}{\snorm{x}}$ from the $\eta$ in \probref{muon} to cancel out the $\snorm{x}$ term in the final line of our derivation of $\snorm{\change y}$. We will cover this scale factor in a dedicated scale factor section below. For now, we will ignore this factor in the naive derivation.  

Perform SVD on the gradient $\grad[\Weights]\lossfunc = U\Sigma V^\top$, then we can solve \probref{muon-final} by:
\begin{align}
    \change \Weights = -\eta \times UV^\top
\end{align}

Now, to not perform SVD every time we try to calculate the update, we can approximate SVD through Newton Schultz iteration, and hence we have \ruleref{base-muon-update}. 

\subsection{Scale Factor}

Using \ruleref{base-muon-update} in SGD Momentum, we can see how Algorithm 1 is constructed. However, Algorithm 1 is not the full implementation in Keller Jordan's GitHub repository \cite{jordan2024muon}. We will need a few more steps. 

Let's first define the RMS norm:
\begin{align}
    \rmsnorm{v} := \sqrt{\frac{1}{d} \sum^{d}_{i=1} v^2_i} = \frac{1}{\sqrt{d}} \times \snorm{v}
\end{align}

Note that in practice, modern architecture often normalize activation using LayerNorm or RMSNorm, which keeps hidden representation close to unit RMS magnitude. Using this fact, we can analyze updates under the assumption:
\begin{align*}
    \rmsnorm{x} \le 1
\end{align*}

Now, looking back at $\frac{1}{\snorm{x}}$ factor from \probref{muon-final}. Once we rewrite the derivation using RMS norm, observe that this factor naturally disappears using our previous assumption. 

Another consequence of the approximately unit RMS magnitude is that when we use RMS norm, the norm stays consistent across dimension, as opposed to spectral norm, where the value is dependent on vector dimension \cite{bernstein2025deriving}. Thus, using RMS norm allows the learning rates to be transferable across dimensions, which is a desirable property. 

Now, we define the RMS-to-RMS operator norm over a weight matrix:
\begin{align}
    \rmsrmsnorm{\Weights} := \max_{x \ne 0} \frac{\rmsnorm{\Weights x}}{\rmsnorm{x}} = \sqrt{\frac{\fanin}{\fanout}} \times \snorm{\Weights}
\end{align}

Replacing the spectral norm in \probref{muon-final}, we can write down the problem again:
\begin{align}
    \min_{\change \Weights} \langle \grad[\Weights]\lossfunc, \change \Weights \rangle \text{ subject to } \snorm{\change \Weights} \le \sqrt{\frac{\fanout}{\fanin}} \eta
\end{align}

We'll also update the solution accordingly:
\begin{align}
    \change \Weights = - \sqrt{\frac{\fanout}{\fanin}}\eta \times UV^\top
\end{align}

We now move on to Conv2D layers. Since Conv2D layers have a 4-dimensional weight space, the previous methods don't apply. However, there is very similar norm we can use, define $k$ to be the kernel size of the layer, and $\Weights \in \R^{\fanout, \fanin, k, k}$ \cite{bernstein2024modulardualitydeeplearning}:
\begin{align}
    k^2 \max^k_{i,j = 1} \rmsrmsnorm{W_{\cdot\cdot ij}}
\end{align}

Using this norm, we can also obtain a similar solution for Conv2D layers:
\begin{align}
    \change \Weights[\cdot\cdot ij] = - k^2\sqrt{\frac{\fanout}{\fanin}}\eta \times U_{ij}V_{ij}^\top
\end{align}

For ease of notation, we will only reference the linear version of the algorithm in the following sections, but this solution replaces the linear solution whenever working with Conv2D layers. 

\subsection{Putting Muon Together}

With our theoretical understanding of Muon, we can now piece together the full implementation of Muon. Our main focus here is line 6 in Algorithm 1, the update rule. In particular, we'll add the scale factor and decoupled weight decay, which is empirically proven to improve Muon's scalability\cite{liu2025muonscalablellmtraining}. Combining these, we have our updated algorithm:

\begin{algorithm}
\caption{Practical Muon Implementation}
\begin{algorithmic}[1]
    \Require Learning rate $\eta$, momentum $\mu$, weight decay $\lambda$
    \State Initialize: $\Momentum[0] \gets 0$
    \For{$t = 1, \dots$}
        \State Compute gradient $\Gr[t] \gets \grad[\Weights]\lossfunc(\Weights[t-1])$
        \State $\Momentum[t] \gets \eta \Momentum[t-1] + \Gr[t]$
        \State $\Orthogonal[t] \gets$ NewtonSchultz5$(\Momentum[t])$
        \State Update parameters $\Weights[t] \gets \Weights[t-1] - \eta(\sqrt{\frac{\fanout}{\fanin}}\Orthogonal[t] + \lambda \Weights[t-1])$
    \EndFor

    \State \Return $\Weights[t]$
\end{algorithmic}
\end{algorithm}

Importantly, Muon should NOT be used in any of the following layers \cite{chollet2015keras}:
\begin{enumerate}
    \item Embedding layer
    \item Final output fully connected layer
    \item Any $\{0,1\}$-D variables
\end{enumerate}

\subsection{More on Scale Factors}

Interestingly, this scale factor we derived above is the same as that of $\mu P$ \cite{yang2022tensorprogramsvtuning}. In addition to this, there are a few more versions of Muon that all differ slightly in scale factor, each with its respective rationales. We will not elaborate too much as their proposers have laid out justification; we will only mention them here to highlight the choice to be made. Below, we will provide the main weight update line, and all else remains the same \cite{kexuefm-11416} \arjun{add some more details, this is really cool (ie compatability with adam), also provide sources for each variation}:

\begin{align}
    \Weights[t] \gets \Weights[t-1] - \eta(\sqrt{\max(1, \frac{\fanout}{\fanin})}\Orthogonal[t] + \lambda \Weights[t-1])
\end{align}
\begin{align}
    \Weights[t] \gets \Weights[t-1] - \eta(0.2 \times \sqrt{\max(\fanout, \fanin)}\Orthogonal[t] + \lambda \Weights[t-1])
\end{align}

\section{Manifold Muon}

Now that we have a theoretical basis for Muon, let's turn our attention to Manifold Muon. The major difference here is that now our weight space is constrained to the Stiefel manifold:
\[\text{Stiefel}(m,n) := \{W \in \mathbb{R}^{m \times n} \mid W^TW = I_n\}\]

Another note on notation, in this section, we will let $\Tangent$ denote the matrix tangent to the stiefel manifold at matrix $\Weights$ and $\dualvar$ the dual variable. 

We will also use $\msign$ \arjun{write out the diagonal version for clarity ie $U \text{sign}(V) U^\top$} to denote the matrix orthogonalization, which, in practical implementation, can either be implemented using the aforementioned Newtown Schultz iteration or the Polar Express algorithm \cite{amsel2025polarexpressoptimalmatrix}.

\subsection{Theoretical Basis}

We will gloss over the formal derivation of the weight update of Manifold Muon, as Jeremy Bernstein has worked out the math already. Instead, we will focus on the aforementioned scale factor. 

Let's look at the original Manifold Muon problem as written \cite{bernstein2025manifolds}:
\begin{align}
    \min_{A \in \mathbb{R}^{m \times n}} \text{trace}(\transpose{\Gr} \Tangent) \quad \text{such that} \quad \snorm{\Tangent} \le \eta  \quad \text{and} \quad \transpose{\Tangent}\Weights + \transpose{\Weights}\Tangent = 0
\end{align}

The first minimization is our problem, the inequality is the update size constraint we mentioned in equation (3), and the final equality is the tangent constraint. 

We can see that the problem used the spectral norm. However, as noted in section 2.2, to make the learning rate transferable, we can switch to using the RMS norm (implication being that we need to satisfy the same "dense" assumption). Doing so would change the inequality to:
\[\rmsrmsnorm{\Tangent} \le \eta\]

Using the definition, we rewrite this to:
\[\snorm{\Tangent} \le \sqrt{\frac{\fanout}{\fanin}}\eta\]

Now we can use the exact same solution, but plugging in $\sqrt{\frac{\fanout}{\fanin}}\eta$ instead of $\eta$, which would allow the learning rate to be transferable. 

So we can take the original solution and write down the algorithm:

\begin{algorithm}[t]
\caption{Vanilla Manifold Muon Update}
\label{alg:manifold_muon}
\begin{algorithmic}[1]
\Require learning rate $\eta$, dual ascent learning rate $\alpha$, iterations $T$, tolerance $\epsilon$

\State $transpose \gets (m < n)$
\If{$transpose$}
    \State $\Weights \gets \transpose{\Weights}$, $\Gr \gets \transpose{\Gr}$
\EndIf

\State $\dualvar \gets -\frac14 (\transpose{\Weights} \Gr + \transpose{\Gr} \Weights)$

\State $\dualvar \gets \textsc{DualAscent}(\Weights, \Gr, \dualvar, \alpha, T, \epsilon)$

\State $\Weights \gets \Weights - \sqrt{\frac{\fanout}{\fanin}}\eta \times \msign[\Gr + 2\Weights\dualvar]$
\State $\Weights \gets \msign[\Weights]$

\If{$transpose$}
    \State $\Weights \gets \transpose{\Weights}$
\EndIf

\State \Return $\Weights$
\end{algorithmic}
\end{algorithm}

\begin{algorithm}[t]
\caption{Dual Ascent (for Vanilla Manifold Muon)}
\label{alg:dual_ascent_manifold_muon}
\begin{algorithmic}[1]
\Require $\Weights, \Gr, \dualvar, \alpha, T, \epsilon$

\For{$t = 1,\dots,T$}
    \State $\Tangent \gets \msign[\Gr + 2\Weights\dualvar]$
    \State $H \gets \transpose{\Weights} \Tangent + \transpose{\Tangent} \Weights$
    \If{$\frobeniusnorm[H]/ \sqrt{|H|} < \epsilon$}
        \State \textbf{break}
    \EndIf
    \State $\dualvar \gets \dualvar - \alpha (1 - t/T) H$
\EndFor

\State \Return $\dualvar$
\end{algorithmic}
\end{algorithm}

\subsection{Aligning with Muon}

TODO: go over how update strategies should follow those suggested in the Muon section (momentum, scale factor, weight decay). This section will wait for more empirical backing

TODO: Weight decay alignment with Muon will have to wait for empirical proof, as the Moonlight article introduces weight decay for a specific reason that we don't have proof for on manifold muon yet

\subsection{ADMM Speedup}

Manifold Muon significantly increases wall clock time per step compared to AdamW \cite{bernstein2025manifolds}, and a major reason is the dual ascent inner loop. We can improve the speed significantly using the ADMM algorithm \cite{buchanan2025mmuonadmm}. Buchanan has a great post explaining the theoretical basis, so we will just write out the algorithm. Notation-wise, we use $\MSubproblem$, $\XSubproblem$, and $\OmSubproblem$ to denote the subproblems in the final solution to solve for $\dualvar$. 

The ADMM algorithm doesn't rely on a break condition like the above dual ascent algorithm. Instead, it takes two hypermeters: $T$ and $\rho$, for iteration count and penalty, respectively. For our purposes, we take Buchanan's conclusions from the parameter sweep and use $(10, 4.0)$ for the two hyperparameters. \footnote{should we do our own experiments or edits to optimize this instead of choosing an arbitrary hyperparameter}

\begin{algorithm}[t]
\caption{ADMM Dual Solve for Manifold Muon (Dual Ascent Section Only)}\footnotemark
\label{alg:manifold_muon_admm_dual}
\begin{algorithmic}[1]
\Require $\Weights, \Gr$, penalty term $\rho$, iterations $T$ 

\State $\dualvar[0] \gets -\frac14 ({\Weights}^\top \Gr + {\Gr}^\top \Weights)$
\State $\XSubproblem[0] \gets \Gr + 2 \Weights \dualvar$
\State $\OmSubproblem[0] \gets \mathbf{0}$

\For{$t = 1,\dots,T$}
    \State $P \gets \transpose{\Weights} (\frac{1}{\rho}\OmSubproblem[t-1] + \XSubproblem[t-1] - \Gr)$
    \State $\dualvar[t]\gets \frac14 (P + \transpose{P})$

    \State $\MSubproblem[t] \gets 2\Weights + \Gr \dualvar[t] - \frac{1}{\rho}\OmSubproblem[t-1]$
    \State $\XSubproblem[t] \gets \frac12 (\MSubproblem[t] - \frac{1}{\rho}\msign[\MSubproblem[t]])(I + \msign[\transpose{\MSubproblem[t]}{\MSubproblem[t]}- \frac{1}{\rho^2}I])$

    \State $\OmSubproblem[t] \gets \OmSubproblem[t-1] + \rho(\XSubproblem[t] - 2\Weights \dualvar[t] - \Gr)$
\EndFor
\State \Return $\dualvar$
\end{algorithmic}
\end{algorithm}
\footnotetext{$P$ used in the algorithm is a temporary variable that don't hold inherent meaning}

\subsection{PDHG Speedup}

TODO: Second speedup method for dual ascent using the PDHG algorithm \cite{cesista2025steepestdescentfinsler} Should run comparison as Buchanan proposed

\clearpage

\section{Experiments}

\arjun{some becnhmarks: cifar10, cifar-100, simple transoformer, tiny qwen, gan}. 

TODO: GAN results, also should run more experiments on Muon and Manifold Muon esp. with different implementation choices. 

Note: Need to choose a good baseline test for basic comparisons, is GAN a good choice here?

List of potential experiments to run with final, corrected algorithms
\begin{enumerate}
    \item GAN tests like before, run w. different implementation
    \item Simpler baseline on CIFAR-10/CIFAR-100 to compare across specific implementations?
    \item ADMM vs. PDHG vs. Basic DA comparison
    \item Potential speed advantage on tasks that want orthogonality/normalization?
    \item This idea is from chatgpt, but maybe transformers?
\end{enumerate}

\newpage

\printbibliography

\end{document}